
\subsection{Resumen: matrices, vectores e inversas del proceso}\label{sec:algebra-proceso}

Los cuatro pasos anteriores se resolvieron con matrices y vectores, y cada vez que hay que deshacer una transformación aparece una
inversa. Esta subsección recoge en un solo lugar lo que ya se calculó en cada paso, con los valores de la foto del casco; la teoría
está en el marco teórico (sección «\nameref{sec:inversa}»).

\paragraph{Matrices del proceso.} El cuadro~\ref{tab:matrices-proceso} muestra las matrices en el orden en que se crean. Los tamaños
cuadran en cada producto: la retroproyección de los $N$ píxeles se hace de una vez con
\begin{equation}
\underbrace{P}_{3\times N}=\underbrace{K^{-1}}_{3\times3}\ \underbrace{H}_{3\times N}\ \underbrace{\operatorname{diag}(Z_1,\dots,Z_N)}_{N\times N},
\end{equation}
donde $(3\times3)(3\times N)$ da $3\times N$ y multiplicar por la diagonal a la derecha escala cada columna. Como en el programa cada punto es una
fila, los vértices quedan en $V\in\mathbb{R}^{N\times3}$ al pasar a los ejes del objeto con la transpuesta de la rotación:
\begin{equation}
V=\bigl(P^{T}-\vect{1}\,\vect{c}^{T}\bigr)\,R_o^{T},\qquad R_o=\operatorname{diag}(1,-1,-1),
\end{equation}
con $\vect{1}$ un vector de unos de longitud $N$ (resta el centro $\vect{c}$ a todas las filas) y $R_o^{T}=R_o$.

\begin{table}[htbp]
\caption{Matrices del proceso (tamaños con la foto del casco)}
\label{tab:matrices-proceso}
\centering\small
\begin{tabular}{@{}p{2.3cm}p{3.1cm}p{4.9cm}p{3.2cm}@{}}
\toprule
Matriz & Tamaño & Qué guarda & Paso \\
\midrule
$I$ & $1280\times960\times3$ & el color de cada píxel & adquisición \\
$M$ & $1280\times960$ & 1 en el objeto, 0 en el fondo & paso 1 \\
$K$ & $3\times3$ & el modelo de la cámara: $f=753{,}449$, $(c_x,c_y)=(480,640)$ & paso 3 \\
$H$ & $3\times32\,657$ & los píxeles del casco como columnas $(u,v,1)^{T}$ & paso 3 \\
$\operatorname{diag}(Z)$ & $32\,657\times32\,657$ & la profundidad de cada píxel, entre $284$ y $381$ mm & paso 3 \\
$P$ & $3\times32\,657$ & los puntos 3D en los ejes de la cámara & paso 3 \\
$R_o$ & $3\times3$ & giro de $180^\circ$ alrededor de $x$ & paso 3 \\
$\mathcal{M}_o$ & $4\times4$ & $R_o$ y la traslación $-R_o\vect{c}$ juntas (coordenadas homogéneas) & paso 3 \\
$V$ & $65\,314\times3$ & los vértices de la malla (frente y atrás) & paso 4 \\
$F$ & $130\,624\times3$ & los índices de los tres vértices de cada triángulo & paso 4 \\
\bottomrule
\end{tabular}
\end{table}

Además de los productos, se usan la \emph{transpuesta} (para pasar de columnas a filas y para invertir rotaciones) y el
\emph{determinante}: el de $K$, $\det K=f^2=567\,685{,}5$, comprueba que se puede invertir, y el de la matriz $3\times3$ de los vértices de
cada triángulo da el volumen del tetraedro que forma con el origen.

\paragraph{Vectores del proceso.} Los datos que viajan por el proceso son vectores, y las tres operaciones del curso (norma, producto punto y
producto cruz) aparecen en lugares concretos (cuadro~\ref{tab:vectores-proceso}).

\begin{table}[htbp]
\caption{Vectores del proceso y operaciones que se les aplican (valores del casco)}
\label{tab:vectores-proceso}
\centering\small
\begin{tabular}{@{}p{2.3cm}p{3.0cm}p{4.9cm}p{4.0cm}@{}}
\toprule
Vector & Qué representa & Valor & Operación \\
\midrule
$\vect{c}_A=(r,g,b)$ & color del píxel A & $(1,\,1,\,1)$ & norma de la diferencia: $\lVert\vect{c}_A-\vect{c}_B\rVert=405{,}4$ \\
$\vect{p}=(u,v,1)^{T}$ & píxel en coordenadas homogéneas & $(488,576,1)^{T}$ & producto $K^{-1}\vect{p}$ \\
$\vect{r}=K^{-1}\vect{p}$ & dirección del rayo del píxel & $(0{,}0106,-0{,}0849,1)^{T}$ & norma $1{,}00366$; ángulo con el eje óptico $4{,}89^\circ$ \\
$\vect{P}=Z\,\vect{r}$ & punto 3D (ejes de la cámara) & $(3{,}19,-25{,}56,300{,}87)^{T}$ mm & se multiplica por $K$ para volver al píxel \\
$\vect{e}_1,\vect{e}_2$ & dos lados de un triángulo & $\lVert\vect{e}_1\rVert=1{,}1899$ & resta de vértices \\
$\vect{n}=\vect{e}_1\times\vect{e}_2$ & normal del triángulo & $(-0{,}1844,-0{,}3248,1{,}3803)^{T}$ & producto cruz; $\hat n_z=0{,}9653>0$ mira al observador \\
\bottomrule
\end{tabular}
\end{table}

El producto punto aparece en el ángulo entre rayos, $\cos\theta=\vect{r}_A\cdot\vect{r}_C/(\lVert\vect{r}_A\rVert\lVert\vect{r}_C\rVert)$, y, en las
habitaciones, en la ecuación del plano $\vect{n}\cdot\vect{P}=d$ y en la intersección rayo--plano $Z=d/(\vect{n}\cdot\vect{r})$.

\paragraph{Matrices de $3\times3$ y de $4\times4$.} Solo las matrices cuadradas se pueden invertir (cuando su determinante no es cero). En el proceso hay
cuatro de $3\times3$ y dos de $4\times4$ (cuadro~\ref{tab:cuadradas}). Se cumple $\det K\cdot\det K^{-1}=f^2\cdot f^{-2}=1$, y $R_o$ y las dos matrices
homogéneas tienen determinante $1$, que es la señal de que no deforman ni cambian volúmenes.

\begin{table}[htbp]
\caption{Matrices cuadradas del proceso, con su determinante y su inversa}
\label{tab:cuadradas}
\centering\small
\begin{tabular}{@{}p{2.2cm}p{1.6cm}p{3.4cm}p{3.5cm}p{2.7cm}@{}}
\toprule
Matriz & Tamaño & Determinante & Inversa & Para qué sirve \\
\midrule
$K$ & $3\times3$ & $f^2=567\,685{,}5$ & $K^{-1}$, $\det K^{-1}=1{,}762\times10^{-6}$ & del píxel al rayo \\
$R_o$ & $3\times3$ & $1$ & $R_o^{-1}=R_o^{T}=R_o$ & girar los ejes \\
$T$ (color) & $3\times3$ & $0{,}2363$ & $T^{-1}$ & de YCbCr a RGB \\
$[\vect{p}_0;\vect{p}_1;\vect{p}_2]$ & $3\times3$ & $111{,}093$ (de un triángulo) & no se invierte & su determinante da el volumen \\
$\mathcal{M}_o$ & $4\times4$ & $1$ & $\mathcal{M}_o^{-1}=\begin{pmatrix}R_o^{T}&\vect{c}\\\vect{0}^{T}&1\end{pmatrix}$ & cambio de ejes completo \\
$\mathcal{M}$ (alto) & $4\times4$ & $1$ & $\mathcal{M}^{-1}=\mathcal{M}^{T}$ & nodo del modelo generativo \\
\bottomrule
\end{tabular}
\end{table}

\paragraph{Inversas del proceso.} Cada vez que se pasa de un lado a otro hay una matriz y su inversa. El cuadro~\ref{tab:inversas-proceso} reúne las de $3\times3$
y $4\times4$ y dice cómo se obtiene cada una; ninguna necesita un método general porque todas tienen una estructura que la simplifica.

\begin{table}[htbp]
\caption{Inversas de las matrices de $3\times3$ y $4\times4$ y cómo se calculan}
\label{tab:inversas-proceso}
\centering\small
\begin{tabular}{@{}p{2.1cm}p{2.9cm}p{4.2cm}p{4.2cm}@{}}
\toprule
Matriz & Inversa & Cómo se calcula & Para qué sirve \\
\midrule
$K$ ($3\times3$, triangular) & $K^{-1}$ & restar $c_x$, $c_y$ veces la fila 3 y dividir entre $f$ (Gauss--Jordan con 3 operaciones) & del píxel al rayo: $\vect{r}=K^{-1}\vect{p}$ \\
$R_o$ ($3\times3$, ortogonal) & $R_o^{-1}=R_o^{T}$ & trasponer, sin cuentas & volver de los ejes del objeto a los de la cámara \\
$T$ ($3\times3$) & $T^{-1}$ & Gauss--Jordan & pasar de YCbCr a RGB \\
$\mathcal{M}_o$ ($4\times4$) & $\mathcal{M}_o^{-1}$ & por bloques: $R_o^{T}$ y $\vect{c}$ & volver del objeto a la cámara con una sola matriz \\
$\mathcal{M}$ ($4\times4$, alto) & $\mathcal{M}^{-1}=\mathcal{M}^{T}$ (rotación pura) & por bloques, $s^{-1}R^{T}$ y $-s^{-1}R^{T}\vect{t}$ & transformación de nodo del modelo generativo \\
\bottomrule
\end{tabular}
\end{table}

También son cuadradas e invertibles la matriz de la DCT del JPEG ($C$, de $8\times8$, con $C^{-1}=C^{T}$) y $\operatorname{diag}(Z)$ (de $N\times N$, con inversa
$\operatorname{diag}(1/Z)$), pero son matrices grandes que se invierten sin hacer cuentas, por eso no están en el cuadro.

\paragraph{La inversa de la cadena completa.} El paso del píxel al vértice es una cadena de transformaciones, y la inversa de una cadena es la cadena de
las inversas en orden contrario, $(AB)^{-1}=B^{-1}A^{-1}$. Para el píxel A, ir del píxel al vértice es $\vect{p}\ \to\ \vect{P}=Z\,K^{-1}\vect{p}\ \to\
\vect{p}_{\text{obj}}=R_o(\vect{P}-\vect{c})$, y volver es
\begin{equation}
\vect{P}=R_o^{T}\,\vect{p}_{\text{obj}}+\vect{c}=R_o^{T}\begin{pmatrix} 6{,}549 \\ -10{,}544 \\ 78{,}664 \end{pmatrix}+\begin{pmatrix} -3{,}354 \\ -36{,}101 \\ 379{,}535 \end{pmatrix}=\begin{pmatrix} 3{,}195 \\ -25{,}557 \\ 300{,}871 \end{pmatrix},
\end{equation}
\begin{equation}
K\vect{P}=\begin{pmatrix} 146\,825{,}0 \\ 173\,301{,}6 \\ 300{,}9 \end{pmatrix}=Z\,\vect{p},
\end{equation}
y al dividir entre $Z=300{,}87$ se recupera el píxel $(488,576,1)$ de partida (con la matriz de $4\times4$ $\mathcal{M}_o$, el mismo viaje de vuelta es un solo
producto, $\mathcal{M}_o^{-1}(\vect{p}_{\text{obj}},1)$): la cadena y su inversa se cancelan, que es lo que
comprueban las pruebas automáticas del servidor ($KK^{-1}=I$ y proyectar un punto retroproyectado devuelve el mismo píxel).

\paragraph{La matriz que no se puede invertir.} La única matriz del proceso sin inversa es la de la proyección, $A=\begin{pmatrix}f&0&c_x-u\\0&f&c_y-v\end{pmatrix}$,
que es de $2\times3$: no es cuadrada, tiene rango $2$ y su espacio nulo es la recta del rayo (sección «\nameref{sec:casco-retro}»). Por eso una sola foto no
alcanza para recuperar la profundidad y hay que agregar información: la silueta y la red de profundidad en los objetos, y los planos del cuarto en las
habitaciones, donde $Z=d/(\vect{n}\cdot\vect{r})$ convierte el sistema en una ecuación con una sola incógnita.
